% Part E fragments owned by cycles agent. No preamble.
\subsection{E.1: the missing saturation condition on 14:9}
\label{sec:effectivity-saturation}

\paragraph{E.1a (Proved: saturation as a precise polarity).}
Let $G/\mathbb Q$ be a reductive algebraic group, with its category of
finite-dimensional rational representations. Let $\mathcal I$ be the
set of complex cocharacters, modulo conjugation if desired. For
$C\subset\mathcal I$ and a class $\mathcal E$ of representations,
define
\[
\begin{split}
 E(C)&=\{V:\text{every weight of every }\nu\in C
                 \text{ on }V_{\mathbb C}\text{ is nonnegative}\},\\
 I(\mathcal E)&=\{\nu:\text{every weight of }\nu
                 \text{ on every }V\in\mathcal E
                 \text{ is nonnegative}\},\\
 \operatorname{Sat}(C)&=I(E(C)).
\end{split}
\]
Then $\operatorname{Sat}$ is extensive, monotone, and idempotent, and
$E(\operatorname{Sat}(C))=E(C)$. The exact missing condition is
\begin{equation}\label{eq:effective-saturation}
 C=I(E(C)).
\end{equation}
The saturated cocharacter sets are in inclusion-reversing bijection
with the full representation subcategories that are intersections of
the $P(\nu)$. Additional positivity, Tate-twist, weight-zero, and
Galois-invariance axioms are imposed on the objects on both sides;
they do not replace \eqref{eq:effective-saturation}.

\emph{Proof.}
The defining condition gives the Galois connection
$\mathcal E\subset E(C)$ if and only if $C\subset I(\mathcal E)$.
In particular $C\subset I(E(C))$. Applying the order-reversing map
$E$ gives $E(I(E(C)))\subset E(C)$, while the reverse inclusion is
true from the definition of $I$. Thus $E(I(E(C)))=E(C)$, which also
proves idempotence. If $\mathcal E=\bigcap_{\nu\in C}P(\nu)=E(C)$,
its full set of effective cocharacters is $I(\mathcal E)$ and is
saturated. Conversely a saturated $C$ is recovered from the
intersection defining $E(C)$. This proves the bijection. The same
argument applies to pro-reductive groups by using their actual
finite-dimensional representations; it does not require choosing
one finite-dimensional maximal torus of a pro-group.

\paragraph{E.1b (Proved lattice form).}
For a torus with character lattice $M$ and cocharacter lattice
$N=\operatorname{Hom}(M,\mathbb Z)$, assume $C\subset N$ is Galois
stable. Then
\[
 C^\vee=\{m\in M:\langle m,c\rangle\ge0\quad\forall c\in C\},
 \qquad
 \operatorname{Sat}(C)=
 \{v\in N:\langle m,v\rangle\ge0\quad\forall m\in C^\vee\}.
\]
To check that rational representations detect every listed character,
we take the representation whose weights are the finite Galois orbit
of $m$. Galois stability of $C$ makes the whole orbit effective
exactly when $m$ is in $C^\vee$. For a split connected reductive group,
the same formula works on a maximal torus for a Weyl-stable cone:
move $m$ into a dominant chamber; the irreducible representation of
highest weight $m$ has all its weights in the convex hull of its Weyl
orbit. The Weyl-stable polar therefore makes that whole representation
effective. Over $\mathbb Q$ we use a rational representation containing
the finite Galois orbit, with multiplicity if descent requires it;
multiplicities do not change inequalities.

\paragraph{E.1c (Refuted automaticity for additive cones).}
We give complete data compatible with the normalization
$\langle\varepsilon,j\rangle=2$. Let
\[
 T=\mathbb G_m\times R^1_{\mathbb Q(i)/\mathbb Q}\mathbb G_m,
 \qquad N=\mathbb Z^2,
 \qquad \overline{(a,b)}=(a,-b),
 \quad \varepsilon=(1,0)\in M,
 \quad j=(2,0)\in N.
\]
Set
\[
 C=\mathbb N(1,1)+\mathbb N(1,-1)
   =\{(a,b)\in\mathbb Z^2:a\ge |b|, a\equiv b\pmod2\}.
\]
It is an additive Galois-stable cone, contains $j$, and has
$\varepsilon$ strictly positive away from the origin.
The character $\varepsilon$ is in the interior of its real polar.
Its integral polar is
$C^\vee=\{(u,v):u\ge|v|\}$, because the two generator inequalities
are $u+v\ge0$ and $u-v\ge0$. Therefore
\[
 \operatorname{Sat}(C)=\{(a,b):a\ge|b|\}\supsetneqq C;
 \qquad (1,0)\in\operatorname{Sat}(C)\setminus C.
\]
The two cones define exactly the same effectivity structure. This
is a counterexample to automatic saturation under the additive-cone
requirements alone.

Even imposing root saturation is not enough. The monoid
\[
 C_{\mathrm{open}}=\{0\}\cup
     \{(a,b)\in\mathbb Z^2:a>|b|\}
\]
is root saturated: $nv\in C_{\mathrm{open}}$ for an integer $n>0$
implies $v\in C_{\mathrm{open}}$. It has the same integral polar and
the same bipolar as $C$, but omits $(1,1)$ and $(1,-1)$. Thus the
required saturation includes closure as detected by integral
half-spaces, not merely divisibility in the cocharacter lattice.

\paragraph{E.1d (Proved automaticity under a stronger meaning of ``cone'').}
There is an important qualification to the preceding answer. Suppose
``cone'' is defined in advance to mean
$C=K\cap N$, where $K\subset N_{\mathbb R}$ is a closed convex cone
and $\varepsilon$ lies in the interior of $K^\vee$.
Then $C$ is saturated as above, in particular when $K$ is rational
polyhedral. The conclusion holds without rational polyhedrality:
a separating
normal for a lattice point outside $K$ can be perturbed into the
interior of $K^\vee$ while keeping its value negative at that point,
then approximated by a rational normal and scaled to an integral
character. This character is nonnegative on $K\cap N$ and excludes
the point from the integral bipolar.

For completeness, the perturbation uses
$m_t=m+t\varepsilon$ for sufficiently small $t>0$, where
$m\in K^\vee$ satisfies $\langle m,v\rangle<0$.
Since $\varepsilon$ is interior, $m_t$ is interior; a sufficiently
close rational point remains there and remains negative on $v$.
The separating-hyperplane theorem supplies $m$ because $K$ is closed
and convex. This proves both inclusions of the integral bipolar.

Thus the manuscript's unspecified quotation-marked ``c\^one'' matters.
Saturation is not a consequence of the stated additive requirements;
it is automatic if the stronger closed-convex-lattice definition is
already part of those requirements. Formula
\eqref{eq:effective-saturation} removes that ambiguity entirely.
It also gives the strongest exact classification requested in E.1.

\paragraph{E.1e (Audit of the portion of R13 used here).}
The closure of $P(\nu)$ under sums, subobjects and tensor products
follows by taking weights, and a positive
$\langle\varepsilon,\nu\rangle$ makes sufficiently many Tate twists
of every fixed representation effective. For a cone with the strict
interior condition, the same twist works uniformly on that cone:
normalize $\langle\varepsilon,c\rangle=1$, use compactness of the
closed section, and bound each of the finitely many weights of the
representation on it. This verifies the step of R13 needed in E.1.
We do not use E.1 as a proof of the motivic conjecture that every
cohomologically effective motive is effective. Lead E1, concerning
the Galois orbit and Frobenius slopes, belongs to E.2 and is not
decided by these purely representation-theoretic statements.

